\section{Distribuciones de Bose--Einstein y Fermi--Dirac}
\label{v:rec:ocupacion:section:02}

La formulación de equilibrio parte del Hamiltoniano y del número total de
ocupaciones,
\begin{equation}
 H=\sum_kE_k^{\rm occ} n_k,
 \qquad
 N=\sum_k n_k,
 \qquad
 K_\mu:=H-\mu N.
 \label{v:rec:eq:hamiltoniano-gran-canonico-rev3}
\end{equation}
Para \(\beta>0\), el estado gran canónico es
\begin{equation}
 \rho_{\beta,\mu}
 =\frac{e^{-\beta K_\mu}}{\mathcal Z_{\beta,\mu}},
 \qquad
 \mathcal Z_{\beta,\mu}
 =\operatorname{Tr}(e^{-\beta K_\mu}).
 \label{v:rec:eq:estado-gran-canonico-rev3}
\end{equation}

\begin{theorem}[Factorización gran canónica y leyes de ocupación]
\label{v:rec:thm:factorizacion-gran-canonica-rev3}
Para niveles de energía \(E_k^{\rm occ}\) con degeneraciones \(g_k\), la
traza fermiónica factoriza como
\begin{equation}
 \mathcal Z_{\rm FD}
 =\prod_k\bigl(1+e^{-\beta(E_k^{\rm occ}-\mu)}\bigr)^{g_k},
 \qquad
 \langle n_k\rangle_{\rm FD}
 =\frac{g_k}{e^{\beta(E_k^{\rm occ}-\mu)}+1}.
 \label{v:rec:eq:factorizacion-fd-rev3}
\end{equation}
En el sector bosónico, bajo la condición
\(\mu<\inf_kE_k^{\rm occ}\),
\begin{equation}
 \mathcal Z_{\rm BE}
 =\prod_k\bigl(1-e^{-\beta(E_k^{\rm occ}-\mu)}\bigr)^{-g_k},
 \qquad
 \langle n_k\rangle_{\rm BE}
 =\frac{g_k}{e^{\beta(E_k^{\rm occ}-\mu)}-1}.
 \label{v:rec:eq:factorizacion-be-rev3}
\end{equation}
\end{theorem}

\begin{proof}
El operador \(K_\mu\) es una suma de operadores de ocupación conmutantes, de
modo que su traza factoriza por modos. En cada modo fermiónico se suman las
ocupaciones \(0\) y \(1\), y el factor correspondiente es
\(1+u_k\), donde
\(u_k=e^{-\beta(E_k^{\rm occ}-\mu)}\). En cada modo bosónico se suman todas
las ocupaciones no negativas y se obtiene la serie geométrica
\((1-u_k)^{-1}\); la condición sobre \(\mu\) garantiza \(u_k<1\). Las
potencias \(g_k\) incorporan la degeneración. Finalmente,
\[
 \langle n_k\rangle
 =-\frac1\beta\frac{\partial}{\partial E_k^{\rm occ}}
   \log\mathcal Z
\]
produce las dos leyes de ocupación
\cite{Bose1924,Einstein1924,Fermi1926,Dirac1926,Pathria2011}.
\end{proof}
